\documentclass[12pt,a4paper]{article}
\usepackage[T2A]{fontenc}
\usepackage[utf8]{inputenc}
\usepackage[russian]{babel}
\usepackage{amsmath,amssymb}
\usepackage{geometry}
\geometry{margin=2.2cm}

\begin{document}

\textbf{Задача.} Привести к нормальной форме $\lambda$-терм
\[
((\lambda a.(\lambda b.\, b\, b)\ (\lambda b.\, b\, b))\ b)\ ((\lambda c.\,(c\, b))\ (\lambda a.\, a)).
\]

Для удобства обозначим
\[
\Omega = (\lambda b.\, b\, b)\ (\lambda b.\, b\, b).
\]

Тогда  терм можно записать так:
\[
((\lambda a.\, \Omega)\ b)\ ((\lambda c.\,(c\, b))\ (\lambda a.\, a)).
\]

Рассматриваем редекс
\[
(\lambda a.\,\Omega)\,b.
\]

Тогда
\[
(\lambda a.\,\Omega)\,b \to_\beta \Omega[a:=b]=\Omega.
\]

Поэтому
\[
((\lambda a.\,\Omega)\,b)\ ((\lambda c.\,(c\, b))\ (\lambda a.\, a))
\to_\beta
\Omega\ ((\lambda c.\,(c\, b))\ (\lambda a.\, a)).
\]

Теперь посмотрим на $\Omega$:
\[
\Omega = (\lambda b.\, b\, b)\ (\lambda b.\, b\, b).
\]

Если выполнить $\beta$-редукцию, получится
\[
(\lambda b.\, b\, b)\ (\lambda b.\, b\, b)
\to_\beta
(\lambda b.\, b\, b)\ (\lambda b.\, b\, b).
\]

То есть снова получается то же самое выражение. Значит,
\[
\Omega \to_\beta \Omega,
\]
и этот процесс будет повторяться бесконечно.


Следовательно, исходный $\lambda$-терм не приводится к нормальной форме.

\textbf{Ответ:} нормальной формы нет.

\end{document}